\documentclass[11pt]{article}
\usepackage[margin=1in]{geometry}
\usepackage{booktabs}
\usepackage{hyperref}

\title{CAI6734 Weekly Report 4: Completing LA-CDM Training}
\author{Omer Kahveci}
\date{October 20, 2026}

\begin{document}
\maketitle

\section{Project Objective}

This week I continued the LA-CDM reproduction using Meta-Llama-3-8B-Instruct
and the four-condition MIMIC clinical decision-making cohort. The immediate
goal was to complete one 300-step training cycle despite the memory limit of
the two NVIDIA L4 GPUs available through the course allocation.

\section{Accomplishments}

The training schedule assigned 100 optimizer steps each to Decision Agent
group-relative policy optimization (GRPO), Hypothesis Agent supervised
fine-tuning, and Hypothesis Agent confidence-calibration GRPO. One L4 GPU
held the trainable model and the other ran vLLM to generate clinical
trajectories. The final job used four generations per case, a forward/backward
microbatch size of one, a 128-token limit per generated completion, and a
2,048-token trajectory limit. The model used LoRA adapters, BF16 computation,
and non-reentrant gradient checkpointing.

To fit the training GPU's memory, I changed the trainer to split each GRPO
group into one-example forward and backward passes, compute reference-model
log probabilities in no-gradient chunks, and calculate only the loss active
in the current training phase. For Hypothesis Agent fine-tuning, I restricted
the scored logits to the short supervised suffix instead of the full prompt.
I also enabled non-reentrant gradient checkpointing and reduced the maximum
trajectory length from 4,096 to 2,048 tokens. Checkpoints every 25 steps
allowed interrupted runs to resume without repeating completed training.

The resumed run completed all 300 optimizer steps and ended with a clean
Slurm finish line on October 20. The final validation pass covered 240 cases
with four rollouts per case, or 960 predictions. Table~\ref{tab:validation}
places these results beside the earlier untrained zero-shot test run, which
used one rollout for each of 240 test cases.

\begin{table}[ht]
  \centering
  \caption{Untrained test results and trained validation results. These use
  different data splits and rollout counts and are descriptive, not a
  controlled before-and-after comparison.}
  \label{tab:validation}
  \begin{tabular}{lrr}
    \toprule
    Metric & Untrained test & Trained validation \\
    \midrule
    Diagnosis accuracy, all predictions & 8.75\% & 56.46\% \\
    Diagnosis accuracy, diagnosed only & 72.41\% & 67.50\% \\
    Undiagnosed fraction & 87.92\% & 16.35\% \\
    Action-format reward & 0.1208 & 0.8365 \\
    Macro F1, all predictions & 15.54\% & 57.59\% \\
    Hypothesis accuracy & 65.25\% & 63.06\% \\
    Expected calibration error & 10.87\% & 7.89\% \\
    \bottomrule
  \end{tabular}
\end{table}

The reported training loss was 0.005625. The final validation took about
2.48 hours. The large difference in all-prediction accuracy is encouraging,
but the different splits and rollout counts prevent attributing that
difference to training alone.

For context, the original LA-CDM paper trained Qwen-2.5-7B-Instruct and
reported the test-set results in Table~\ref{tab:paper}. The paper trained
until validation convergence, taking approximately three days; our Llama run
completed one 300-step cycle. The different model, training duration, and
evaluation setup make these published values reference points rather than
direct targets for our validation run
(\href{https://arxiv.org/abs/2506.13474}{Bani-Harouni et al., 2026},
Table~1, Section~6, and Appendix~C).

\begin{table}[ht]
  \centering
  \caption{Published trained Qwen LA-CDM test results.}
  \label{tab:paper}
  \begin{tabular}{lr}
    \toprule
    Metric & Paper result \\
    \midrule
    Mean class accuracy & 81.3\% \\
    Micro F1 & 84.1\% \\
    Macro F1 & 81.3\% \\
    Average test cost & \$1,295.61 \\
    Hypothesis accuracy & 81.9\% \\
    Expected calibration error & 3.7\% \\
    \bottomrule
  \end{tabular}
\end{table}

\section{Challenges}

Earlier jobs exceeded the training GPU's roughly 22~GB usable memory on long
trajectories, including an out-of-memory failure during backpropagation at
step 90. Once the memory changes allowed longer runs, the Hypothesis Agent
phase exposed a separate batching error: its target rows did not match the
four GRPO rollout rows. I corrected its microbatch indexing and weighted each
microbatch by its share of supervised target tokens. The completed run passed
both loss-phase transitions.

Some trajectories still contained invalid action formats or unknown test
names; the 16.35\% undiagnosed fraction shows that completing the action
protocol remains a meaningful limitation.

\section{Next Steps}

This result establishes that the full three-phase pipeline can train and
resume on the course GPU allocation. The side-by-side figures do not yet
establish improvement against the zero-shot baseline under matched evaluation
conditions. I will verify the saved final adapter and
evaluate it on the fixed 240-case test set using the baseline evaluation
settings. I will then compare accuracy, undiagnosed fraction, format reward,
test use, and calibration, and inspect representative invalid-action examples.

\section{AI Disclaimer}

I used OpenAI Codex to help interpret logs, revise training code, and draft
this report. I reviewed the technical claims and remain responsible for the
final report.

\end{document}
